\clearpage
\VerdierPhysicalPageAnchor{0171}%
\begingroup\small\setlength{\parskip}{0pt}\setlength{\emergencystretch}{1.5em}
\VerdierSemanticAnchor{ch4-4-3-3}
\medskip\noindent\textbf{4.3.3.}\quad
我们特别关注各项
$E\!\left[\begin{smallmatrix}n\\p-r+1,p,p+1,p+r\end{smallmatrix}\right]$，
$p\in\mathbb Z$，$r\in\mathbb Z$，$r>0$，因为在适当的假设下，可以由
这些项确定 $H^n(-\infty,+\infty)$ 的关联分次对象。为了回到
\emph{同上引文}的记号，我们令：

\VerdierSemanticAnchor{ch4-4-3-3-1}
\[
\left\{
\begin{aligned}
E\!\left[\begin{smallmatrix}n\\p-r,p-1,p,p+r-1\end{smallmatrix}\right]
  &=\mathrm I_r^{-p,n+p},\\
&\textit{即}:\\[-2pt]
\mathrm I_r^{p,q}
  &=E\!\left[\begin{smallmatrix}p+q\\-p-r,-p-1,-p,-p+r-1\end{smallmatrix}\right],
  &&p,q\in\mathbb Z,\quad 1\leq r\leq+\infty,\\
&\text{相应的微分记为：}\\[-2pt]
d_r\colon\mathrm I_r^{p,q}&\longrightarrow
  \mathrm I_r^{p+r,q-r+1}.
\end{aligned}
\right.
\tag{4.3.3.1}
\]

我们将经常使用如下记号：

\VerdierSemanticAnchor{ch4-4-3-3-2}
\[
\left\{
\begin{aligned}
E\!\left[\begin{smallmatrix}n\\p-r+1,p,p+1,p+r\end{smallmatrix}\right]
  &=\mathrm{II}_{r+1}^{n-p,p},\\
&\textit{即}:\\[-2pt]
\mathrm{II}_r^{p,q}
  &=E\!\left[\begin{smallmatrix}p+q\\q-r+2,q,q+1,q+r-1\end{smallmatrix}\right],
  &&p,q\in\mathbb Z,\quad 2\leq r\leq+\infty,\\
&\text{相应的微分记为：}\\[-2pt]
d_r\colon\mathrm{II}_r^{p,q}&\longrightarrow
  \mathrm{II}_r^{p+r,q-r+1}.
\end{aligned}
\right.
\tag{4.3.3.2}
\]

我们将在每种情形中指明所使用的记号。注意，使用记号 (4.3.3.1) 时，
$\mathrm I_r^{p,q}$ 对 $r\geq1$ 有定义；而使用记号 (4.3.3.2) 时，
$\mathrm{II}_r^{p,q}$ 只对 $r\geq2$ 有定义。

\VerdierSemanticAnchor{ch4-4-3-3-spectral-sequence}
各项
$E\!\left[\begin{smallmatrix}n\\p-r+1,p,p+1,p+r\end{smallmatrix}\right]$，
$n,p\in\mathbb Z$，$1\leq r\leq+\infty$，连同相应的微分，称为
\emph{谱对象} $(H^n(p,q),\delta(p,q,r))$ \emph{所关联的谱序列}。
\endgroup
